\documentclass[11pt]{article}

% ---------- Packages ----------
\usepackage[margin=1in]{geometry}
\usepackage{amsmath, amssymb, amsthm}
\usepackage{mathtools}
\usepackage{enumitem}
\usepackage{graphicx}
\usepackage{booktabs}
\usepackage{array}
\usepackage{hyperref}
\usepackage{xcolor}
\usepackage{fancyhdr}
\usepackage{lastpage}
\usepackage{listings}
\usepackage{courier}

% ---------- Hyperref ----------
\hypersetup{
    colorlinks=true,
    linkcolor=blue,
    urlcolor=blue
}

% ---------- Header/Footer ----------
\pagestyle{fancy}
\fancyhf{}

\lhead{MIT 6.1200J}
\chead{Problem Set 1}
\rhead{Nicolas Iskandar}

\cfoot{\thepage\ /\ \pageref{LastPage}}

% ---------- Custom Commands ----------
\newcommand{\R}{\mathbb{R}}
\newcommand{\N}{\mathbb{N}}
\newcommand{\Z}{\mathbb{Z}}

\newcommand{\answer}{
    \vspace{0.5em}
    \noindent\textbf{Answer.}
}

% ---------- Code Style ----------
\lstset{
    basicstyle=\ttfamily\small,
    breaklines=true,
    frame=single,
    columns=fullflexible,
    keepspaces=true
}

% ---------- Title ----------
\title{
    \vspace{-2em}
    MIT 6.1200J \\
    Problem Set 1
}
\author{Nicolas Iskandar}
\date{}

% ==================================================
\begin{document}

\maketitle

\section*{Problem 2}

\begin{enumerate}[label=\textbf{(\alph*)}]

\item Use a truth table to verify that the propositional formula $(P \land \neg Q)\lor(P \land Q) \equiv P$ is equivalent to P.

\answer

\[
\begin{array}{|c|c|c|c|c|c|}
\hline
P & Q & \neg Q & P \land \neg Q & P \land Q & (P \land \neg Q)\lor(P \land Q) \\
\hline
T & T & F & F & T & T \\
T & F & T & T & F & T \\
F & T & F & F & F & F \\
F & F & T & F & F & F \\
\hline
\end{array}
\]

\[
(P \land \neg Q)\lor(P \land Q) \equiv P
\]

\item Prove that $A = (A - B) \cup (A \cap B)$
for all sets A, B, by showing $x \in A \iff x \in (A - B) \cup (A \cap B)$ for all elements x. How does this relate to part (a)?

\answer

Let \(x\) be arbitrary. Then,
\[
x \in (A - B) \cup (A \cap B)
\]
\[
\iff x \in (A - B) \ \lor\ x \in (A \cap B)
\]
\[
\iff (x \in A \land x \notin B)\ \lor\ (x \in A \land x \in B)
\]
\[
\iff x \in A \land (\neg(x \in B) \lor x \in B)
\]
\[
\iff x \in A \land \text{True}
\]
\[
\iff x \in A
\]

Hence,
\[
x \in (A - B) \cup (A \cap B) \iff x \in A
\]

Therefore,
\[
A = (A - B) \cup (A \cap B)
\]

\textbf{Relation to part (a):} This identity is an instance of the propositional equivalence
\[
(P \land \neg Q)\ \lor\ (P \land Q) \equiv P
\]
where \(P\) corresponds to \(x \in A\) and \(Q\) corresponds to \(x \in B\).

\end{enumerate}

\section*{Problem 3}

\begin{enumerate}[label=\textbf{(\alph*)}]

\item The predicate $\text{sixish}(n)$ is defined as ``n is a natural number whose first digit is 6.''

\answer
\[
\text{sixish}(n) := \exists k. \Big( 6 \cdot 10^k \leq n \wedge n < 7 \cdot 10^k \Big)
\]

\item ``There are exactly two natural numbers $n$ such that $Q(n)$ is true.''

\answer
\[
\exists x \; \exists y. \Big( x \neq y \wedge Q(x) \wedge Q(y) \wedge \forall z. \big( Q(z) \implies (z = x \vee z = y) \big) \Big)
\]

\item Goldbach's Conjecture: ``Every even integer greater than 2 can be written as the sum of two primes.''

\answer
\[
\forall n. \Big( \big( n > 2 \wedge \exists k.\, n = 2k \big) \implies \exists p \; \exists q. \big( \text{isPrime}(p) \wedge \text{isPrime}(q) \wedge p + q = n \big) \Big)
\]

\item ``Every prime larger than 100 is sixish.''

\answer
\[
\forall n. \Big( \big( \text{isPrime}(n) \wedge n > 100 \big) \implies \text{sixish}(n) \Big)
\]

\item Is the predicate $R(n) := \exists b.\, n = 2b + 5$ equivalent to $\text{isOdd}(n)$?

\answer

No, they are not equivalent. We show $R(n) \not\iff \text{isOdd}(n)$ by counterexample.

\textbf{Counterexample:} Take $n = 1$.
\begin{itemize}
    \item $\text{isOdd}(1)$ is \textbf{true}: set $m = 0$, then $2(0) + 1 = 1$. \checkmark
    \item $R(1)$ is \textbf{false}: $R(1)$ requires $\exists b \in \mathbb{N}$ such that $2b + 5 = 1$,
    i.e., $2b = -4$. Since $b \in \mathbb{N}$, no such $b$ exists.
\end{itemize}

Therefore the predicates are \textbf{not equivalent}.
\end{enumerate}

\newpage

\section*{Problem 4}

\begin{enumerate}[label=\textbf{(\alph*)}]

\item We prove the following statement by induction on $n$.

\textbf{Inductive Hypothesis} $P(n)$: For any $n$ numbers $x_1, x_2, \ldots, x_n \in \mathbb{R}$, if $Q(x_1 \cdot x_2 \cdots x_n)$ is true, then at least one of $Q(x_1), Q(x_2), \ldots, Q(x_n)$ is true.

\textbf{Base Case} ($n = 1$): Suppose $Q(x_1)$ is true. Then trivially, $Q(x_1)$ is true. So $P(1)$ holds.

\textbf{Inductive Step}: Assume $P(n)$ holds for some $n \geq 1$ (inductive hypothesis). We want to show $P(n+1)$: that for any $x_1, \ldots, x_{n+1} \in \mathbb{R}$, if $Q(x_1 \cdot x_2 \cdots x_{n+1})$ is true, then at least one of $Q(x_1), \ldots, Q(x_{n+1})$ is true.

\textbf{Assumption}: Suppose $Q(x_1 \cdot x_2 \cdots x_{n+1})$ is true.

\textbf{Goal}: Show at least one of $Q(x_1), \ldots, Q(x_{n+1})$ is true.

Let $y = x_1 \cdot x_2 \cdots x_n$, so that
\[
    x_1 \cdot x_2 \cdots x_{n+1} = y \cdot x_{n+1}.
\]
Since $Q(y \cdot x_{n+1})$ is true, the given property of $Q$ tells us that at least one of $Q(y)$ or $Q(x_{n+1})$ is true. We consider two cases:

\begin{itemize}
    \item \textbf{Case 1}: $Q(x_{n+1})$ is true. Then we are done.
    \item \textbf{Case 2}: $Q(y) = Q(x_1 \cdot x_2 \cdots x_n)$ is true. By the inductive hypothesis $P(n)$, at least one of $Q(x_1), Q(x_2), \ldots, Q(x_n)$ is true.
\end{itemize}

In either case, at least one of $Q(x_1), \ldots, Q(x_{n+1})$ is true. This completes the inductive step.

By induction, $P(n)$ holds for all integers $n \geq 1$. $\blacksquare$

\item We choose $Q(x) :=$ ``$x$ is irrational.''

With this choice, part (a) states: if $x_1 \cdot x_2 \cdots x_n$ is irrational, then at least one of $x_1, \ldots, x_n$ is irrational, which is exactly Theorem 1.

However, to apply part (a), we must first verify that $Q$ satisfies the required property from part (a). This gives us the following lemma:

\textbf{Lemma}: For any $x, y \in \mathbb{R}$, if $x \cdot y$ is irrational, then at least one of $x$ or $y$ is irrational.

\item \textbf{Proof of Lemma} (by contradiction):

Suppose for contradiction that $x \cdot y$ is irrational, but neither $x$ nor $y$ is irrational. Then both $x$ and $y$ are rational, so we may write
\[
    x = \frac{a}{b}, \qquad y = \frac{c}{d},
\]
for some integers $a, b, c, d$ with $b \neq 0$ and $d \neq 0$. Then
\[
    x \cdot y = \frac{a}{b} \cdot \frac{c}{d} = \frac{ac}{bd}.
\]
Since $a, b, c, d$ are integers, $ac$ and $bd$ are integers, and $bd \neq 0$. Therefore $x \cdot y = \frac{ac}{bd}$ is rational, so a contradiction, since we assumed $x \cdot y$ is irrational. $\blacksquare$

\newpage

\textbf{Proof of Theorem 1}:

Let $Q(x) :=$ ``$x$ is irrational.'' The Lemma establishes that $Q$ satisfies the hypothesis of part (a): for any $x, y \in \mathbb{R}$, if $Q(x \cdot y)$ is true, then at least one of $Q(x)$ or $Q(y)$ is true. Therefore, by part (a), for any $n \geq 1$ and any $x_1, \ldots, x_n \in \mathbb{R}$, if $Q(x_1 \cdot x_2 \cdots x_n)$ is true (i.e., the product is irrational), then at least one of $Q(x_1), \ldots, Q(x_n)$ is true (i.e., at least one of the numbers is irrational). This proves Theorem 1. $\blacksquare$

\end{enumerate}

\end{document}